\chapter{My Journey at UFV}
\setcounter{section}{0}

\section{The PFG as the Culmination of My University Journey}

When I arrived at Universidad Francisco de Vitoria (UFV), I was someone who had just finished
secondary school with a clear interest in computer science but a very narrow view of what the field
could become. I had always enjoyed working with computers, taking them apart, learning how software
worked and trying to build things on my own. There was no family pressure behind the choice;
informatics had been a personal vocation from early on. At that point, my main ambition was to
become a software developer at one of the large technology companies. I imagined a career almost
entirely defined by development.

The choice of UFV was guided by recommendation from people I trusted, and once I attended the first
sessions I felt that the academic environment and the human focus of the institution matched what I
was looking for. I did not yet fully understand the value of the humanistic approach that the UFV
adds to its engineering programmes; that understanding came progressively, over the years.

The first two years of the degree, I was not the most focused version of myself. I tended to leave
things until the last moment, trusted that effort under pressure would be enough and approached the
work without a real sense of direction. The change came towards the end of the second year and the
beginning of the third year. Something shifted in how I understood the relationship between effort,
time and results. From that point on, I stopped treating the degree as a sequence of assignments to
complete and started treating it as a long-term project to invest in. I became more organised,
fought harder for my own objectives and started to take responsibility for the depth of what I was
learning, not only for the grade. I am still far from perfect at this, but I can clearly see the
distance between who I was at the start and who I am now.

Two people stand out as decisive influences during this journey. The first is Mois\'es Mart\'inez
Mu\~n\'oz, supervisor of this project, whose way of thinking about Artificial Intelligence and
research has been an inspiration throughout the final years of the degree. The second is Olga
Pe\~nalba, director of the Escuela Polit\'ecnica Superior, who taught the Projects courses in which
I was particularly engaged, and through which I obtained a Matr\'icula de Honor in one of the
iterations. Both of them, in different ways, helped me see the engineering profession as something
broader and more demanding than the narrow developer career I had pictured at the beginning.

The choice of this specific PFG did not come from a personal connection with the deaf community or
with accessibility as a previous interest. It came from a curiosity about computer vision, an
attraction to Artificial Intelligence as a discipline and the conviction that the final project of
the degree should be ambitious enough to push me into a depth I had not reached yet. The project
emerged from my own reflection on what kind of work I wanted to do, not from an externally proposed
topic. Choosing sign language recognition gave the project a real and meaningful application, while
keeping the technical depth I was looking for. In that sense, the PFG reflects two properties of
who I am: a genuine interest in Artificial Intelligence as a field of study, and a willingness to
commit fully to an objective once I decide it is the right one.

\section{Connection with My Professional Future}

I am currently doing a software engineering internship at Fintonic in parallel with the development
of this project. The experience has been positive and I have learned a lot, but it has also
clarified something important: development is enjoyable and I am comfortable doing it, yet it is
not where I want my career to settle. The PFG has confirmed that my real interest lies in
Artificial Intelligence, in deep learning specifically, and in the engineering practice that makes
machine learning systems reliable, reproducible and deployable.

The project has taught me things that neither the degree alone nor the internship alone could have
taught me. The first lesson is patience. I had never worked on something this long and this
intense, and there were moments of real frustration when progress felt invisible despite hours of
work. The second lesson is project management at a scale I had not experienced before. Maintaining
a coherent direction across months of experimentation, keeping the code reproducible and the
results auditable, and recovering from methodological pivots without losing the thread, are skills
I did not have when I started. I believe the final outcome reflects that growth.

Looking ahead, my objective is to pursue a Master's degree in Artificial Intelligence and, ideally,
a Doctorate, preferably outside Spain. If a doctoral programme aligned with my interests does not
become available, I will look for a role strongly focused on Artificial Intelligence rather than on
general software development. The areas I most want to deepen are deep learning and MLOps, because
they combine the modelling work I find intellectually rewarding with the engineering rigour needed
to bring those models into reliable use. One discovery from this project that I did not anticipate
is that I would not mind being a researcher, and I find the idea of teaching at university level
genuinely appealing. Both are paths I want to keep open as I continue.

The humanistic dimension of the UFV has played a role here that I did not fully see while it was
happening. The way the university trains engineers to think about their work, beyond the purely
technical layer, gives an additional perspective that I now consider valuable. It is part of what
allows me to see Artificial Intelligence not only as a set of algorithms to optimise, but as a
discipline that affects how people communicate, how they are represented in datasets and how they
are served by the systems we build. That perspective is visible in the ethical analysis of this
project and will remain visible in whatever I do next.

Throughout all of this, the people who have most influenced who I am professionally and personally
are my family and the professors who shaped my view of the field. The PFG closes one stage of that
journey and opens a new one, in which I want to keep growing, keep specialising in what I love and,
in time, contribute something meaningful to the field of Artificial Intelligence.